Large language models (LLMs) are increasingly used as software-engineering assistants, yet it remains difficult to tell when a correct answer reflects faithful program reasoning. We study this question through generated Java tasks that isolate two forms of program comprehension: control-flow reachability, where a model must decide whether one method can call another directly or indirectly, and data-flow tracing, where the same benchmark design can require following values across calls, parameters, and local state. The benchmark exposes graph family, context size, source-target depth, syntactic distractors, and requested reasoning budget as explicit experimental factors, and it evaluates both final Boolean answers and the traces that justify them. In the current control-flow study, three reasoning-capable models are evaluated on a common slice of 108,000 samples over contexts of 100, 250, and 500 methods, depths 1--25, and requested budgets from 512 to 16,384 tokens. The results show that final-answer accuracy substantially overstates reasoning quality: Gemma reaches 90.1\% final accuracy but only 57.2\% exact-chain accuracy, GPT-OSS reaches 79.0\% and 36.3\%, and Qwen reaches 89.1\% and 66.4\%. Budget helps, but only model-specifically: Qwen converts budget into large final and chain gains, Gemma saturates early, and GPT-OSS improves final answers faster than trace fidelity. These findings suggest that reasoning-budget evaluations for software engineering should report path-level evidence, not only final labels.
